\section{Operator System}

The evolution of the triadic fields is written in terms of a minimal operator family acting on
\(\phi(\mathbf{x}, t)\), \(\vec{V}(\mathbf{x}, t)\), and \(R(\mathbf{x}, t)\). For clarity, we write the dynamics in the form
\begin{align}
    \partial_t \phi &= \mathcal{D}_\phi[\phi] + \mathcal{A}_\phi[\phi, \vec{V}, R] + \mathcal{C}_\phi[\phi, \vec{V}, R] + \mathcal{N}_\phi[\phi, R] + \mathcal{S}_\phi[\phi], \\
    \partial_t \vec{V} &= \mathcal{D}_V[\vec{V}] + \mathcal{A}_V[\phi, \vec{V}, R] + \mathcal{C}_V[\phi, \vec{V}, R] + \mathcal{N}_V[\vec{V}, R] + \mathcal{S}_V[\vec{V}], \\
    \partial_t R &= \mathcal{D}_R[R] + \mathcal{A}_R[\phi, \vec{V}, R] + \mathcal{C}_R[\phi, \vec{V}] + \mathcal{N}_R[R, \phi, \vec{V}] + \mathcal{S}_R[R].
\end{align}

\subsection{Diffusion Operator}

We use standard Laplacian diffusion with possibly distinct coefficients:
\begin{align}
    \mathcal{D}_\phi[\phi] &= D_\phi \nabla^2 \phi, \\
    \mathcal{D}_V[\vec{V}] &= D_V \nabla^2 \vec{V}, \\
    \mathcal{D}_R[R] &= D_R \nabla^2 R,
\end{align}
where \(D_\phi, D_V, D_R > 0\) are diffusion constants.

\subsection{Alignment Operator}

Alignment terms drive local coherence between the fields. A simple form is
\begin{align}
    \mathcal{A}_\phi[\phi, \vec{V}, R] &= \alpha_\phi\, R \left( \phi^\ast - \phi \right), \\
    \mathcal{A}_V[\phi, \vec{V}, R] &= \alpha_V\, R \left( \vec{V}_\parallel - \vec{V} \right), \\
    \mathcal{A}_R[\phi, \vec{V}, R] &= \alpha_R \left( C(\phi, \vec{V}) - R \right),
\end{align}
where \(\phi^\ast\) is a preferred scalar profile, \(\vec{V}_\parallel\) is the component of \(\vec{V}\) aligned with a local direction (e.g., a gradient or external field), and \(C(\phi, \vec{V})\) is a scalar coherence functional (e.g., a normalized correlation between \(\nabla \phi\) and \(\vec{V}\)). The coefficients \(\alpha_\phi, \alpha_V, \alpha_R\) control alignment strength.

\subsection{Coupling Operator}

Coupling terms mediate cross-field influence:
\begin{align}
    \mathcal{C}_\phi[\phi, \vec{V}, R] &= \beta_\phi\, \nabla \cdot (R \vec{V}), \\
    \mathcal{C}_V[\phi, \vec{V}, R] &= \beta_V\, R \nabla \phi, \\
    \mathcal{C}_R[\phi, \vec{V}] &= \beta_R\, \left( \vec{V} \cdot \nabla \phi \right),
\end{align}
with \(\beta_\phi, \beta_V, \beta_R\) setting the strength of scalar–vector–resonance coupling.

\subsection{Activation Operator}

Nonlinear activation terms allow thresholded or saturating behavior. A generic choice is
\begin{align}
    \mathcal{N}_\phi[\phi, R] &= \gamma_\phi\, f_\phi(\phi, R), \\
    \mathcal{N}_V[\vec{V}, R] &= \gamma_V\, f_V(\vec{V}, R), \\
    \mathcal{N}_R[R, \phi, \vec{V}] &= \gamma_R\, f_R(R, \phi, \vec{V}),
\end{align}
where \(f_\phi, f_V, f_R\) are smooth nonlinear functions (e.g., sigmoidal, polynomial, or piecewise-defined) and \(\gamma_\phi, \gamma_V, \gamma_R\) are activation coefficients. In applications, these can encode resonance thresholds, saturation, or gain control.

\subsection{Stabilization Operator}

Stabilization terms enforce bounded dynamics and suppress runaway growth:
\begin{align}
    \mathcal{S}_\phi[\phi] &= -\lambda_\phi\, \phi, \\
    \mathcal{S}_V[\vec{V}] &= -\lambda_V\, \vec{V}, \\
    \mathcal{S}_R[R] &= -\lambda_R\, R,
\end{align}
with \(\lambda_\phi, \lambda_V, \lambda_R > 0\) providing linear damping. Additional higher-order stabilizers (e.g., \(-\eta_R R^3\)) can be introduced if needed.
